% Options for packages loaded elsewhere
\PassOptionsToPackage{unicode}{hyperref}
\PassOptionsToPackage{hyphens}{url}
%
\documentclass[
  11pt,
  a4paper,
]{article}
\usepackage{amsmath,amssymb}
\usepackage{iftex}
\ifPDFTeX
  \usepackage[T1]{fontenc}
  \usepackage[utf8]{inputenc}
  \usepackage{textcomp} % provide euro and other symbols
\else % if luatex or xetex
  \usepackage{unicode-math} % this also loads fontspec
  \defaultfontfeatures{Scale=MatchLowercase}
  \defaultfontfeatures[\rmfamily]{Ligatures=TeX,Scale=1}
\fi
\usepackage{lmodern}
\ifPDFTeX\else
  % xetex/luatex font selection
  \setmainfont[]{DejaVu Serif}
  \setsansfont[]{DejaVu Sans}
  \setmonofont[]{DejaVu Sans Mono}
\fi
% Use upquote if available, for straight quotes in verbatim environments
\IfFileExists{upquote.sty}{\usepackage{upquote}}{}
\IfFileExists{microtype.sty}{% use microtype if available
  \usepackage[]{microtype}
  \UseMicrotypeSet[protrusion]{basicmath} % disable protrusion for tt fonts
}{}
\makeatletter
\@ifundefined{KOMAClassName}{% if non-KOMA class
  \IfFileExists{parskip.sty}{%
    \usepackage{parskip}
  }{% else
    \setlength{\parindent}{0pt}
    \setlength{\parskip}{6pt plus 2pt minus 1pt}}
}{% if KOMA class
  \KOMAoptions{parskip=half}}
\makeatother
\usepackage{xcolor}
\usepackage[margin=22mm]{geometry}
\usepackage{longtable,booktabs,array}
\usepackage{calc} % for calculating minipage widths
% Correct order of tables after \paragraph or \subparagraph
\usepackage{etoolbox}
\makeatletter
\patchcmd\longtable{\par}{\if@noskipsec\mbox{}\fi\par}{}{}
\makeatother
% Allow footnotes in longtable head/foot
\IfFileExists{footnotehyper.sty}{\usepackage{footnotehyper}}{\usepackage{footnote}}
\makesavenoteenv{longtable}
\setlength{\emergencystretch}{3em} % prevent overfull lines
\providecommand{\tightlist}{%
  \setlength{\itemsep}{0pt}\setlength{\parskip}{0pt}}
\setcounter{secnumdepth}{5}

\usepackage{amsmath,amssymb,booktabs,longtable,array,etoolbox,xurl,fvextra,graphicx}
\usepackage{seqsplit}
\let\singularitytexttt\texttt
\renewcommand{\texttt}[1]{\singularitytexttt{\seqsplit{#1}}}
\graphicspath{{docs/}}
\usepackage{fancyhdr}
\pagestyle{fancy}
\fancyhf{}
\fancyhead[L]{\small Singularity: Relational Heights and Arithmetic Visibility}
\fancyfoot[C]{\thepage}
\setlength{\headheight}{15pt}
\setlength{\emergencystretch}{3em}
\setlength{\tabcolsep}{3pt}
\renewcommand{\arraystretch}{1.12}
\AtBeginEnvironment{longtable}{\footnotesize\everymath{\scriptstyle}}
\fvset{breaklines=true,breakanywhere=true}
\allowdisplaybreaks
\makeatletter
\renewcommand*\l@section{\@dottedtocline{1}{1.5em}{3.6em}}
\renewcommand*\l@subsection{\@dottedtocline{2}{3.8em}{4em}}
\renewcommand*\@pnumwidth{2.5em}
\renewcommand*\@tocrmarg{3.5em}
\makeatother
\AtBeginDocument{\hypersetup{colorlinks=true,linkcolor=teal,urlcolor=teal}}
\ifLuaTeX
  \usepackage{selnolig}  % disable illegal ligatures
\fi
\IfFileExists{bookmark.sty}{\usepackage{bookmark}}{\usepackage{hyperref}}
\IfFileExists{xurl.sty}{\usepackage{xurl}}{} % add URL line breaks if available
\urlstyle{same}
\hypersetup{
  pdftitle={Exact Observation Depth for Integer-Height Graph Populations},
  pdfauthor={Singularity project contributors},
  hidelinks,
  pdfcreator={LaTeX via pandoc}}

\title{Exact Observation Depth for Integer-Height Graph Populations}
\usepackage{etoolbox}
\makeatletter
\providecommand{\subtitle}[1]{% add subtitle to \maketitle
  \apptocmd{\@title}{\par {\large #1 \par}}{}{}
}
\makeatother
\subtitle{Connected supports, delayed distinctions, and sharp
reconstruction}
\author{Singularity project contributors}
\date{18 October 2026}

\begin{document}
\maketitle

\hypertarget{abstract}{%
\section*{Abstract}\label{abstract}}
\addcontentsline{toc}{section}{Abstract}

For a finite simple graph, we count nonnegative integer vertex heights
jointly by their sum and a pinned quadratic edge energy. We prove a
connected-support expansion for the formal logarithm: the coefficients
through total height R depend only on connected induced vertex sets of
size at most R, decorated with their ambient degrees. This gives an
exact bounded-degree compiler whose graph-size cost is linear at fixed
R. The expansion identifies the first coefficient at which equal-order
regular graphs of girth at least g can differ: it is the difference of
their g-cycle counts times an explicit polynomial. Our central
reconstruction theorem determines every n-vertex simple 2-regular graph
from the pinned series through volume floor(n/2), at every fixed real
energy weight strictly between zero and one. The budget is sharp for
every n at least six. In particular, a cycle of length 2g and two
disjoint g-cycles agree below volume g and have a nonzero explicit
discrepancy at g. Their dominating-terminal blocks remain
indistinguishable below that volume under every fixed closing context.
We supply complete formal-series proofs, independent finite enumeration
receipts, cycle reconstruction for all 57 simple 2-regular graphs on at
most fourteen vertices, and retained bounded-degree timing comparisons.
The connected-cluster principle and low-order graph-coefficient methods
have established antecedents; the claim here is the stated
height-population specialization and its exact reconstruction and
observation-depth theorems.

\hypertarget{population-budget-and-contribution}{%
\section{Population, budget and
contribution}\label{population-budget-and-contribution}}

Let H=(V,E) be a finite simple loopless graph. Heights are functions h:V
to the nonnegative integers. They remain attached to the original
labeled vertices; no quotient by graph automorphisms is taken. Define \[
J(h)=\sum_{v\in V}h_v,\qquad
E_H(h)=\sum_{uv\in E}(h_u-h_v)^2,
\] and, for an integer \(\lambda\ge0\), \[
F_\lambda(H;x,q)=\sum_h x^{J(h)}q^{E_H(h)+\lambda\sum_vh_v^2}.
\tag{1}
\] The coefficient of \(x^j\) is a polynomial in q with nonnegative
integer coefficients. It counts a finite population of
\(\binom{j+n-1}{n-1}\) profiles when \(n=|V|>0\). All identities below
are initially identities in \(\mathbb Q[q][[x]]\); evaluation at
\(q_0\in(0,1)\) is coefficientwise. No infinite analytic expansion or
convergence assertion is needed. We write \(A_H=F_2(H)\) and
\(L_H=F_0(H)\).

A budget R means that coefficients of x through degree R are available.
It bounds total integer height, rather than the number of occupied
vertices alone or the maximum height. A profile at volume R occupies at
most R vertices. The logarithm strengthens this elementary observation:
only connected occupied footprints need to be compiled independently. We
exploit that identity to establish a sharp observation threshold for an
infinite graph class.

The broad computational idea is established. Patel and Regts {[}PR2017,
Section 3{]} compute low-order graph-polynomial information from
connected induced subgraph counts, including colored variants. Their
additivity criterion explains the vanishing of disconnected terms. Their
later induced-count algorithm {[}PR2019{]} supplies further
bounded-degree context. Our ambient-degree labels and unbounded-height
formal series require the specific specialization proved below; we do
not claim a new connected-cluster principle. Bauerschmidt, Park and
Rodriguez {[}BPR2024{]} study the integer-valued Gaussian model on
lattices and its high-temperature scaling, a different question from
finite-volume exact graph reconstruction. Chertkov and Chernyak
{[}CC2006{]} express partition functions through belief-propagation loop
corrections. The present paper uses a direct formal logarithm and does
not claim to originate loop expansions or to provide a new
belief-propagation approximation.

The principal result is Theorem 6: the exact and sharp budget
\(\lfloor n/2\rfloor\) for simple 2-regular interiors of known order
n.~The preceding local identities provide its proof and an
implementation. General completeness of the single fixed-pinning series
\(A_H\) is not asserted.

\hypertarget{exact-connected-support-expansion}{%
\section{Exact connected-support
expansion}\label{exact-connected-support-expansion}}

Expand the exponent in (1) as \[
\sum_v w_vh_v^2-2\sum_{uv\in E}h_uh_v,
\qquad w_v=d_H(v)+\lambda.
\tag{2}
\] For \(T\subseteq V\), let \(Z_T\) be (1) with heights outside T
forced to zero. It is essential that (2) retains the ambient degree
\(d_H(v)\), rather than replacing it by the degree in \(H[T]\): an edge
to a vertex forced to height zero still contributes energy. Set
\(Z_\varnothing=1\), and define \[
\Phi_S=\sum_{T\subseteq S}(-1)^{|S|-|T|}\log Z_T.
\tag{3}
\] The logarithm exists formally because each series has constant
coefficient one.

\textbf{Theorem 1 (connected-support locality; SNG-F55).} For every
nonempty S, \(\Phi_S\) vanishes when \(H[S]\) is disconnected, has no
term below degree \(|S|\), and depends only on \(H[S]\) decorated by
\(w_v\). Moreover, \[
\log F_\lambda(H)=\sum_{\substack{\varnothing\ne S\subseteq V\\H[S]\text{ connected}}}\Phi_S.
\tag{4}
\] Thus the coefficients through degree R are determined by the
multiplicities of decorated connected induced subgraphs with at most R
vertices.

\textbf{Proof.} Give each vertex its own variable \(y_v\), replacing
\(x^{J(h)}\) by \(\prod_vy_v^{h_v}\), and call the resulting series F.
Forcing heights outside T to zero is the substitution \(y_v=0\) outside
T. The substitution commutes with formal logarithms. A monomial of
\(\log F\) supported on U therefore survives in \(\log Z_T\) exactly
when \(U\subseteq T\). Its multiplicity in (3) is \[
\sum_{U\subseteq T\subseteq S}(-1)^{|S|-|T|},
\] which equals one when U=S and zero otherwise. Hence \(\Phi_S\)
consists exactly of monomials with support S. Such monomials have total
degree at least \(|S|\). If \(H[S]\) splits into nonempty components,
the restricted energy is a sum over those components, \(Z_S\) is a
product, and its logarithm is a sum. It has no monomial meeting distinct
components, so the exact-support part vanishes. Finally, the restricted
exponents involve only edges inside S and the labels \(w_v\).
Partitioning all monomials by support proves (4), and specializing every
\(y_v=x\) proves the assertion. \(\square\)

This is a sufficient-statistics theorem, not an inverse equivalence:
equal decorated small-subgraph counts imply equal truncated series.
Equal truncated series need not determine all those counts individually,
since their polynomial contributions can be dependent.

\textbf{Proposition 2 (exact recovery and fixed-budget cost).} Write
\(F_\lambda=\sum_{j\ge0}a_jx^j\) and \(B_r=r[x^r]\log F_\lambda\). Then
\[
a_0=1,\qquad
j a_j=\sum_{r=1}^j B_r a_{j-r},\qquad B_r\in\mathbb Z[q].
\tag{5}
\] For fixed R, maximum degree \(\Delta\) and \(\lambda\), an exact
implementation of (3)--(5) has cost linear in n in an
arithmetic-operation model. Its dependence on R and degree can be
exponential; this is not a large-budget polynomial-time claim or a bound
on coefficient bit cost.

\textbf{Proof.} The identity \(xF'=F\,x(\log F)'\) gives (5). The
inverse of a series in \(1+x\mathbb Z[q][[x]]\) has integer-polynomial
coefficients, so \(xF'/F\) also does, proving the assertion for \(B_r\).
A connected s-vertex set has a spanning tree and a walk of length
\(2(s-1)\) visiting all its vertices. Choosing a root and such a walk
bounds the number of sets by \(n\max(1,\Delta)^{2(s-1)}\). For each set
of size at most R, evaluating all restrictions in (3) takes at most
\(2^R\) restricted calculations. A restriction on t vertices contains at
most \(\binom{R+t}{t}\le4^R\) profiles of total height at most R. Energy
exponents are bounded by \((\Delta+\lambda)R^2\). Truncated polynomial
arithmetic, logarithms and (5) have cost polynomial in this exponent
range and R per elementary profile step. One may therefore use the
conservative bound \[
n\,[8\max(1,\Delta)^2]^R\operatorname{poly}(R,\Delta,\lambda)
\tag{6}
\] in coefficient operations. This bound requires no canonical labeling
or graph-isomorphism oracle: omitting memoization preserves it. Exact
memoization changes constants and sharing, not correctness. \(\square\)

The implementation uses tree encodings for common decorated trees, exact
small canonical forms when inexpensive, and a vertex-set key when
canonicalization would be costly. The fallback loses reuse only. It
never identifies different decorated shapes. This matters because
factorial canonicalization can cost more than the profile arithmetic it
is intended to save.

\hypertarget{the-first-cycle-visible-to-a-budget}{%
\section{The first cycle visible to a
budget}\label{the-first-cycle-visible-to-a-budget}}

For a regular graph of degree d, put \[
a=q^{d+\lambda},\qquad f=q^{-2}-1,
\qquad \kappa_{d,\lambda}=af
=q^{d+\lambda-2}(1-q^2).
\tag{7}
\] The temporary Laurent notation f is harmless: each final expression
below is a polynomial when d is at least two. The girth of an acyclic
graph is understood to be infinite. Let \(N_g(H)\) count unoriented
simple g-cycles.

\textbf{Theorem 3 (first-cycle coefficient; SNG-F56).} Let g be at least
three. Suppose H and H' are d-regular graphs of the same order, both of
girth at least g. Then \[
F_\lambda(H)\equiv F_\lambda(H')\pmod{x^g},
\qquad
[x^g](F_\lambda(H)-F_\lambda(H'))
=(N_g(H)-N_g(H'))\kappa_{d,\lambda}^{\,g}.
\tag{8}
\]

\textbf{Proof.} We first justify the tree counts used in the locality
formula. For a tree T on k vertices, the number of all graph
homomorphisms into a d-regular n-vertex graph is \(nd^{k-1}\), obtained
by choosing the root image and then one of d neighbors for each child.
Group these maps by their equality partition of the vertices. After
identifying equal vertices and suppressing parallel edges, the quotient
is connected. A quotient with a loop has no map into a simple graph. For
k at most g, any nontrivial quotient has fewer than g vertices; an
injective realization of a cyclic quotient would violate the girth
condition. Every contributing proper quotient is therefore a smaller
tree. Induction on k now shows that the number of injective
homomorphisms of T depends only on n,d and T for k at most g.

For k below g an injective tree image has no extra edge among its
vertices: adding one to the tree would create a cycle on at most k
vertices. Thus the number of induced copies of every such T is its
injective homomorphism count divided by \(|\operatorname{Aut}T|\),
independently of H. All ambient-degree labels also agree. Theorem 1
proves equality of logarithms, and hence series, below degree g.

At degree g, contributions from supports smaller than g still agree. An
induced graph on exactly g vertices is either a tree or a chordless
g-cycle whenever it is connected. An injective tree image can have an
extra edge only if its tree is the path on g vertices: the cycle formed
by that edge must use every vertex. Each g-cycle has exactly g spanning
paths. Consequently the induced-path count is the universal injective
path-image count minus \(gN_g(H)\); every other g-vertex tree count is
universal.

Only height one at every support vertex contributes at total degree g.
The coefficient of \(\prod_{v\in S}y_v\) in \(\log F\) is \[
a^g\sum_{\substack{D\subseteq E(H[S])\\(S,D)\text{ connected}}}f^{|D|}.
\tag{9}
\] To verify (9), expand the binary weight on S as
\(a^{|S|}\prod_{e\in E(H[S])}(1+f)\). Partition each chosen edge set
into its connected components. The exponential formula, or direct
expansion of the formal logarithm followed by this component partition,
leaves precisely connected spanning edge sets. A path contributes
\(a^gf^{g-1}\). A cycle contributes \(a^g(gf^{g-1}+f^g)\), from its g
spanning paths and its full edge set. Subtracting the g missing induced
paths leaves \(a^gf^g\) per cycle. This proves the logarithmic
discrepancy. Since all lower coefficients agree and the constant
coefficients are one, the first discrepancy of the series equals that of
their logarithms. Equation (8) follows. \(\square\)

For the cube and eight-vertex Möbius ladder, d=3, g=4 and the square
counts are six and four. Equation (8) gives \(2q^4(1-q^2)^4\) for L and
\(2q^{12}(1-q^2)^4\) for A. The former is exactly the retained reference
discrepancy. The triangular prism and \(K_{3,3}\) give the analogous
volume-three difference from their two versus zero triangles. These
applications explain particular fixtures; they do not assert that cycle
counts determine the entire height series.

\hypertarget{arbitrarily-late-collisions-and-closing-contexts}{%
\section{Arbitrarily late collisions and closing
contexts}\label{arbitrarily-late-collisions-and-closing-contexts}}

\textbf{Corollary 4 (cycle collision family).} For every integer
\(g\ge3\), \[
A_{C_g\sqcup C_g}-A_{C_{2g}}
=2\kappa^g x^g+O(x^{g+1}),
\qquad \kappa=q^2(1-q^2).
\tag{10}
\] In particular, the series agree below volume g and differ at g for
every \(q\in(0,1)\).

\textbf{Proof.} Both graphs have order 2g, degree two and girth at least
g; their g-cycle counts are two and zero. Apply Theorem 3 with
\(\lambda=2\). \(\square\)

There is an independent transfer proof, useful for checking the leading
term. At a fixed cutoff R, restrict indices a,b to \(0,\ldots,R\) in \[
T_{ab}=x^bq^{(a-b)^2+2b^2}.
\] A cyclic matrix product multiplies exactly one x-factor and one
pinning factor for each vertex, so
\(A_{C_m}=\operatorname{tr}(T^m)\bmod x^{R+1}\). The finite-dimensional
trace identity gives \[
(\operatorname{tr}T^g)^2-\operatorname{tr}T^{2g}
=2\operatorname{tr}((\wedge^2T)^g).
\tag{11}
\] The column indexed by c\textless d in the exterior-square matrix
carries \(x^{c+d}\). Minimum degree one occurs only at \((0,1)\), whose
diagonal entry is \[
x\det\begin{pmatrix}1&q^3\\q&q^2\end{pmatrix}
=xq^2(1-q^2).
\] Every closed product of g entries has degree at least g; equality
forces every intermediate pair to be \((0,1)\). Thus (11) gives exactly
the leading term (10). All matrices were finite, and enlarging R
stabilizes every coefficient under discussion.

To state the context consequence precisely, form \(B_H\) by adding
distinct nonadjacent terminals s,t, each joined to every vertex of H. A
closing context C is a fixed finite simple graph with two designated
ordered terminals. It is glued to \(B_H\) by identifying corresponding
terminals and retaining all other vertices separately. Terminal edge
conventions are unchanged by H, and the context has no additional
identification with the interior. Write \(Z_{C,H}\) for the resulting
unpinned local height series.

\textbf{Proposition 5 (context-independent first discrepancy;
SNG-F53--SNG-F54 specialized).} For interiors of the same order m, \[
Z_{C,H}(x,q)=\sum_{j\ge0} A_{H,j}(q)x^j W_{C,j}(x,q),
\qquad W_{C,j}(0,q)=1.
\tag{12}
\] Hence a first difference at interior volume g remains the same
coefficient at volume g under every such closing context. In particular
the family (10) cannot be distinguished earlier by any fixed context.

\textbf{Proof.} Assign heights s,t to the two terminals. For an interior
profile of volume j, the edges from terminals to H contribute \[
m(s^2+t^2)+2\sum_{v\in H}h_v^2-2j(s+t).
\] The first and last terms depend only on m,j,s,t, and the middle term
is precisely the pinning in \(A_{H,j}\). Summing terminal and
context-interior heights defines \(W_{C,j}\); it is a formal series with
coefficients in \(\mathbb Q(q)\). Its constant term comes uniquely from
zero heights on both terminals and all context-interior vertices, and
equals one. Thus (12) is a unit lower-triangular transformation of the
A-sequence, common to all interiors of order m. It preserves the first
nonzero difference. \(\square\)

For g=6, the two dominating blocks have fourteen vertices. Closing each
with a single shared neighbor gives fifteen-vertex graphs. Direct
weak-composition enumeration of these actual glued graphs gives equality
through volume five and, at volume six and \(q=1/2\), \[
2(\tfrac14\cdot\tfrac34)^6
=\frac{729}{8\,388\,608}.
\] This is a finite independent check of the structural identity, not
evidence that the complete series are equal: (10) proves they are
unequal.

\hypertarget{sharp-reconstruction-of-every-2-regular-interior}{%
\section{Sharp reconstruction of every 2-regular
interior}\label{sharp-reconstruction-of-every-2-regular-interior}}

For k at least one, let \(\varphi_k\) be the exact-support contribution
(3) for a k-vertex path, with every vertex assigned ambient degree two
and pinning two. Theorem 1 gives \(\varphi_k=O(x^k)\), so \[
b(x,q)=\sum_{k\ge1}\varphi_k(x,q),\qquad
D_\ell=\log A_{C_\ell}-\ell b
\tag{13}
\] are well-defined formal series, computable to every finite budget.

\textbf{Theorem 6 (sharp 2-regular reconstruction; SNG-F57).} Fix an
order n and a real \(q_0\in(0,1)\). The coefficients of \(A_H(x,q_0)\)
through degree \(R=\lfloor n/2\rfloor\) determine the isomorphism type
of every simple 2-regular graph H of order n.~For every n at least six
this uniform budget is sharp. Over \(\mathbb Q(q)\), or after the same
specialization, the cycle counts are recovered recursively by \[
\log A_H-nb=\sum_{\ell\ge3}c_\ell D_\ell,
\qquad
c_r=\kappa^{-r}[x^r]\left(\log A_H-nb-\sum_{3\le\ell<r}c_\ell D_\ell\right).
\tag{14}
\]

\textbf{Proof.} A simple 2-regular graph is a disjoint union of cycles
of lengths at least three. For a cycle of length \(\ell\), each
connected proper induced subset is a path. There are \(\ell\) copies of
each path size \(1\le k<\ell\). Therefore Theorem 1 gives \[
\log A_{C_\ell}=\ell\sum_{k<\ell}\varphi_k+\Phi_{C_\ell},
\quad
D_\ell=\Phi_{C_\ell}-\ell\sum_{k\ge\ell}\varphi_k.
\] In particular \(D_\ell=O(x^\ell)\). The multilinear path contribution
at size \(\ell\) is \(q^{4\ell}f^{\ell-1}\); the cycle contribution is
\(q^{4\ell}(\ell f^{\ell-1}+f^\ell)\). Their difference shows \[
D_\ell=\kappa^\ell x^\ell+O(x^{\ell+1}).
\tag{15}
\] The series A multiplies over connected components, and logarithms
add. Subtracting the total vertex background \(nb\) proves the first
identity in (14). Equation (15) makes that identity triangular, proving
its recursive inverse. At \(q_0\in(0,1)\), \(\kappa(q_0)\ne0\), so every
division is valid.

Recover all cycle counts with length at most R. Two remaining cycles
would each have length greater than R, requiring at least \(2R+2>n\)
vertices. Hence there is at most one remaining cycle, and its length is
the remaining vertex count. The recovered multiset of lengths determines
H up to isomorphism.

For sharpness, put \(g=\lfloor n/2\rfloor\). If n=2g, compare \(C_{2g}\)
and \(C_g\sqcup C_g\). If n=2g+1, compare \(C_{2g+1}\) and
\(C_g\sqcup C_{g+1}\). For n at least six all component lengths are
admissible. Their logarithmic differences are respectively
\(2D_g-D_{2g}\) and \(D_g+D_{g+1}-D_{2g+1}\), whose first nonzero
coefficients occur at g by (15). Thus no smaller budget distinguishes
every pair. For n=3,4,5 there is only one simple 2-regular isomorphism
type, so no analogous lower bound is claimed. \(\square\)

The hypothesis that n is known is part of the theorem. It is natural for
fixed-order graph comparison and avoids hiding order recovery inside the
final remainder step. The rational implementation takes exact input
coefficients, validates the recovered nonnegative integer multiplicities
and remaining length, and can reject data that are inconsistent with a
2-regular graph.

\hypertarget{reproduction-and-measured-cost}{%
\section{Reproduction and measured
cost}\label{reproduction-and-measured-cost}}

The package uses Python 3.10 or later and only the standard library for
its exact core. From the repository root,
\texttt{PYTHONPATH=src\ python\ scripts/locality\_family.py\ -\/-check}
reproduces the retained structural receipt in
\texttt{results/locality\_family.json}. The receipt compares the
compiler with an independently implemented weak-composition enumerator
on ten explicit or seeded graphs for both \(\lambda=0\) and \(2\),
checks (8) against the prism and cube examples, verifies (10) for
\(3\le g\le8\) with independent enumeration through g=6, enumerates the
fifteen-vertex glued example directly, and reconstructs all 57 cycle
partitions of total order at most fourteen. The sharpness pairs for
every order six through fourteen are retained too.

The manuscript's theorems are proved above in ordinary mathematics. The
connected-support code and these general theorems are not formalized in
Lean. Separate Lean artifacts for the first paper's fixed reference
coefficients do not enlarge the formal-verification scope of this paper.
Equality to an independent enumeration oracle is finite evidence about
the implementation, not a machine-checked proof of its general
correctness.

The following timings are retained wall-clock measurements, not portable
performance guarantees. Graphs are connected cubic graphs produced by
the seeded configuration-model generator in
\texttt{scripts/locality\_family.py}; n is the seed. The separator
baseline uses the same total-volume and height cutoffs, a stated
min-degree elimination order, and exact coefficient comparison wherever
it is run.

\begin{longtable}[]{@{}
  >{\raggedleft\arraybackslash}p{(\columnwidth - 8\tabcolsep) * \real{0.2000}}
  >{\raggedleft\arraybackslash}p{(\columnwidth - 8\tabcolsep) * \real{0.2000}}
  >{\raggedleft\arraybackslash}p{(\columnwidth - 8\tabcolsep) * \real{0.2000}}
  >{\raggedleft\arraybackslash}p{(\columnwidth - 8\tabcolsep) * \real{0.2000}}
  >{\raggedleft\arraybackslash}p{(\columnwidth - 8\tabcolsep) * \real{0.2000}}@{}}
\toprule\noalign{}
\begin{minipage}[b]{\linewidth}\raggedleft
n
\end{minipage} & \begin{minipage}[b]{\linewidth}\raggedleft
Connected supports, R=4
\end{minipage} & \begin{minipage}[b]{\linewidth}\raggedleft
Separator engine, R=4
\end{minipage} & \begin{minipage}[b]{\linewidth}\raggedleft
Elimination width
\end{minipage} & \begin{minipage}[b]{\linewidth}\raggedleft
Connected supports, R=6
\end{minipage} \\
\midrule\noalign{}
\endhead
\bottomrule\noalign{}
\endlastfoot
16 & 0.027 s & 0.243 s & 4 & 0.900 s \\
32 & 0.025 s & 52.611 s & 7 & 0.818 s \\
64 & 0.037 s & not run & 13 & 0.901 s \\
128 & 0.042 s & not run & 27 & 0.796 s \\
\end{longtable}

At n=64 and 128 the R=4 baseline would allocate, over all bags under its
current representation, approximately \(8.19\times10^9\) and
\(4.66\times10^{19}\) height assignments. These are representation-size
estimates, not measured runtimes or lower bounds for every
tree-decomposition algorithm. The R=6 figures have still larger
estimates and were not run. Observed reuse depends on how many decorated
types occur; at R=6 these four graphs contain 32,26,29 and 17 types.
This explains why the finite timings need not increase monotonically
with n.~The theorem guarantees linear graph-size dependence at fixed
parameters even without reuse, rather than a sublinear guarantee.

To rerun timings use
\texttt{PYTHONPATH=src\ python\ scripts/locality\_family.py\ -\/-benchmark};
this updates the retained timing section and is deliberately separate
from deterministic \texttt{-\/-check}. For the accompanying reference
coefficients use
\texttt{PYTHONPATH=src\ python\ scripts/reproduce.py\ -\/-check}. The
full archive includes sources, tests, exact JSON receipts, the
bibliography and rendered manuscripts. The detailed claim register
identifies every finite boundary and the editorial release check.

\hypertarget{scope-and-next-questions}{%
\section{Scope and next questions}\label{scope-and-next-questions}}

Finite agreement below volume g is now explained and attained by an
explicit family, and within the cycle class a sharp inverse theorem is
available. It does not follow that the fixed-pinning series reconstructs
every finite graph. Nonisomorphic graphs with equal complete A-series
remain a legitimate target. The connected-support identity supplies a
sufficient condition for truncated equality, but its aggregated
q-polynomials can merge different decorated statistics. An exhaustive
prefix search therefore cannot be promoted to a completeness theorem.

Past the first cycle, equal cycle counts need not imply equal
coefficients: overlapping cycles and other connected footprints
contribute. The project has additional loop and overlap results
(SNG-F60--SNG-F61) and a complete moment enrichment (SNG-F58). They are
deliberately outside this paper's central proof and are not
prerequisites for Theorem 6. A further algorithmic question is to
combine separators inside the connected footprints with the fixed-budget
compiler, extending the useful budget without requiring bounded
treewidth of the full graph. These questions retain the same explicit
volume, parameter and graph-class boundaries as the proved results.

\hypertarget{references}{%
\section*{References}\label{references}}
\addcontentsline{toc}{section}{References}

{[}BPR2024{]} R. Bauerschmidt, J. Park and P.-F. Rodriguez, \emph{The
Discrete Gaussian model, I. Renormalisation group flow at high
temperature}, Annals of Probability 52 (2024), no. 4, 1253--1359. DOI:
\href{https://doi.org/10.1214/23-AOP1658}{10.1214/23-AOP1658}. Author
manuscript: \href{https://arxiv.org/abs/2202.02286}{arXiv:2202.02286}.

{[}CC2006{]} M. Chertkov and V. Y. Chernyak, \emph{Loop series for
discrete statistical models on graphs}, Journal of Statistical
Mechanics: Theory and Experiment (2006), P06009. DOI:
\href{https://doi.org/10.1088/1742-5468/2006/06/P06009}{10.1088/1742-5468/2006/06/P06009}.
Author manuscript:
\href{https://arxiv.org/abs/cond-mat/0603189}{arXiv:cond-mat/0603189}.

{[}PR2017{]} V. Patel and G. Regts, \emph{Deterministic polynomial-time
approximation algorithms for partition functions and graph polynomials},
SIAM Journal on Computing 46 (2017), no. 6, 1893--1919. DOI:
\href{https://doi.org/10.1137/16M1101003}{10.1137/16M1101003}. Author
manuscript: \href{https://arxiv.org/abs/1607.01167}{arXiv:1607.01167},
Section 3, particularly Lemma 3.5 and the colored version in Theorem
3.7.

{[}PR2019{]} V. Patel and G. Regts, \emph{Computing the number of
induced copies of a fixed graph in a bounded degree graph}, Algorithmica
81 (2019), 1844--1858. DOI:
\href{https://doi.org/10.1007/s00453-018-0511-9}{10.1007/s00453-018-0511-9}.
Author manuscript:
\href{https://arxiv.org/abs/1707.05186}{arXiv:1707.05186}.

\end{document}
